% GENERATED FILE. Edit canonical lessons, then run npm run generate:books.
% canonical-source-hash: fnv1a64:0271ac9c20b63810
% canonical-lessons: UR-C488-valima, UR-C488-susral, UR-C488-bahu, UR-C488-damad, UR-C488-byahna

\chapter{After the Wedding}
\label{ch:ur-after-the-wedding}
\hlchaptermodality{488}

\begin{chapteropening}
\textbf{By the end of this chapter:} \emph{I can talk about the family a marriage brings.}

Complete the last lesson of chapter 488: I can talk about the family a marriage brings.
\end{chapteropening}

\section[\texorpdfstring{\ur{ولیمہ} (valīmā)}{ولیمہ (valīmā)}]{\ur{ولیمہ} (valīmā) — a wedding feast (given by the groom)}

\label{lesson:UR-C488-valima}



\begin{quote}
Before the new one: say the Urdu for a bridegroom's veil of flowers, then the Urdu for to see off.
\end{quote}

\subsection*{You'll want to know: \ur{ولیمہ}}
\textbf{\ur{ولیمہ}} — \emph{valīmā} — \textquotedblleft{}a wedding feast (given by the groom)\textquotedblright{}.

\begin{quote}
\ur{ولیمے} \ur{میں} \ur{بریانی} \ur{تھی۔}
\end{quote}

\begin{quote}
\emph{valīme meṅ biryānī thī.}
\end{quote}

\begin{quote}
There was biryani at the wedding feast.
\end{quote}

Masculine.

\begin{cousinweb}
\textbf{\ur{ولیمہ}} (\emph{valīmā}) is usually traced, through Persian, to Arabic \emph{walīma}, \textquotedblleft{}a wedding feast\textquotedblright{}.
\end{cousinweb}

\begin{grammarlens}[title={What you've built}]
One more word for describing what you lived and felt.
\end{grammarlens}

\subsection*{Your turn}
\begin{itemize}
\raggedright
  \item \emph{Say it:} \emph{valīmā}
  \item \emph{Say it:} \emph{valīmā}, once more
  \item \emph{Say it:} \emph{valīme meṅ biryānī thī.} — There was biryani at the wedding feast.
  \item \emph{Recall:} say the Urdu for a bridegroom's veil of flowers, then the Urdu for to see off, then say \emph{valīmā} again
\end{itemize}

\begin{tcolorbox}[breakable,colback=teal!4,colframe=teal!35!black,title={Before you move on}]
What does \ur{ولیمہ} mean? (\textquotedblleft{}A wedding feast (given by the groom)\textquotedblright{}.) Say it once more.
\end{tcolorbox}
\section[\texorpdfstring{\ur{سسرال} (susrāl)}{سسرال (susrāl)}]{\ur{سسرال} (susrāl) — the in-laws' home}

\label{lesson:UR-C488-susral}



\begin{quote}
Before the new one: say the Urdu for to see off, then the Urdu for a wedding feast.
\end{quote}

\subsection*{You'll want to know: \ur{سسرال}}
\textbf{\ur{سسرال}} — \emph{susrāl} — \textquotedblleft{}the in-laws' home\textquotedblright{}.

\begin{quote}
\ur{شادی} \ur{کے} \ur{بعد} \ur{وہ} \ur{سسرال} \ur{گئی۔}
\end{quote}

\begin{quote}
\emph{shādī ke bād voh susrāl gaī.}
\end{quote}

\begin{quote}
After the wedding she went to her in-laws' home.
\end{quote}

Feminine.

\begin{cousinweb}
\textbf{\ur{سسرال}} (\emph{susrāl}) is inherited from \emph{śvaśuraśālā}, \textquotedblleft{}the father-in-law's house\textquotedblright{}, a form CDIAL reconstructs (CDIAL 12755); \textbf{\ur{سسر}}, \textquotedblleft{}a father-in-law\textquotedblright{}, is its first part.
\end{cousinweb}

\begin{grammarlens}[title={What you've built}]
One more word for describing what you lived and felt.
\end{grammarlens}

\subsection*{Your turn}
\begin{itemize}
\raggedright
  \item \emph{Say it:} \emph{susrāl}
  \item \emph{Say it:} \emph{susrāl}, once more
  \item \emph{Say it:} \emph{shādī ke bād voh susrāl gaī.} — After the wedding she went to her in-laws' home.
  \item \emph{Recall:} say the Urdu for to see off, then the Urdu for a wedding feast, then say \emph{susrāl} again
\end{itemize}

\begin{tcolorbox}[breakable,colback=teal!4,colframe=teal!35!black,title={Before you move on}]
What does \ur{سسرال} mean? (\textquotedblleft{}The in-laws' home\textquotedblright{}.) Say it once more.
\end{tcolorbox}
\section[\texorpdfstring{\ur{بہو} (bahū)}{بہو (bahū)}]{\ur{بہو} (bahū) — a daughter-in-law}

\label{lesson:UR-C488-bahu}



\begin{quote}
Before the new one: say the Urdu for a wedding feast, then the Urdu for the in-laws' home.
\end{quote}

\subsection*{You'll want to know: \ur{بہو}}
\textbf{\ur{بہو}} — \emph{bahū} — \textquotedblleft{}a daughter-in-law\textquotedblright{}.

\begin{quote}
\ur{ہماری} \ur{بہو} \ur{ڈاکٹر} \ur{ہے۔}
\end{quote}

\begin{quote}
\emph{hamārī bahū ḍākṭar hai.}
\end{quote}

\begin{quote}
Our daughter-in-law is a doctor.
\end{quote}

Feminine.

\begin{cousinweb}
\textbf{\ur{بہو}} (\emph{bahū}) is inherited from Sanskrit \emph{vadhū}, \textquotedblleft{}bride, young wife\textquotedblright{}, already \textquotedblleft{}son's wife\textquotedblright{} in later texts (CDIAL 11250).
\end{cousinweb}

\begin{grammarlens}[title={What you've built}]
One more word for describing what you lived and felt.
\end{grammarlens}

\subsection*{Your turn}
\begin{itemize}
\raggedright
  \item \emph{Say it:} \emph{bahū}
  \item \emph{Say it:} \emph{bahū}, once more
  \item \emph{Say it:} \emph{hamārī bahū ḍākṭar hai.} — Our daughter-in-law is a doctor.
  \item \emph{Recall:} say the Urdu for a wedding feast, then the Urdu for the in-laws' home, then say \emph{bahū} again
\end{itemize}

\begin{tcolorbox}[breakable,colback=teal!4,colframe=teal!35!black,title={Before you move on}]
What does \ur{بہو} mean? (\textquotedblleft{}A daughter-in-law\textquotedblright{}.) Say it once more.
\end{tcolorbox}
\section[\texorpdfstring{\ur{داماد} (dāmād)}{داماد (dāmād)}]{\ur{داماد} (dāmād) — a son-in-law}

\label{lesson:UR-C488-damad}



\begin{quote}
Before the new one: say the Urdu for the in-laws' home, then the Urdu for a daughter-in-law.
\end{quote}

\subsection*{You'll want to know: \ur{داماد}}
\textbf{\ur{داماد}} — \emph{dāmād} — \textquotedblleft{}a son-in-law\textquotedblright{}.

\begin{quote}
\ur{میرے} \ur{ابو} \ur{نانی} \ur{کے} \ur{داماد} \ur{ہیں۔}
\end{quote}

\begin{quote}
\emph{mere abbū nānī ke dāmād haiṅ.}
\end{quote}

\begin{quote}
My father is my grandmother's son-in-law.
\end{quote}

Masculine.

\begin{cousinweb}
\textbf{\ur{داماد}} (\emph{dāmād}) is usually traced to Persian \emph{dāmād}, \textquotedblleft{}son-in-law\textquotedblright{}.
\end{cousinweb}

\begin{grammarlens}[title={What you've built}]
One more word for describing what you lived and felt.
\end{grammarlens}

\subsection*{Your turn}
\begin{itemize}
\raggedright
  \item \emph{Say it:} \emph{dāmād}
  \item \emph{Say it:} \emph{dāmād}, once more
  \item \emph{Say it:} \emph{mere abbū nānī ke dāmād haiṅ.} — My father is my grandmother's son-in-law.
  \item \emph{Recall:} say the Urdu for the in-laws' home, then the Urdu for a daughter-in-law, then say \emph{dāmād} again
\end{itemize}

\begin{tcolorbox}[breakable,colback=teal!4,colframe=teal!35!black,title={Before you move on}]
What does \ur{داماد} mean? (\textquotedblleft{}A son-in-law\textquotedblright{}.) Say it once more.
\end{tcolorbox}
\section[\texorpdfstring{\ur{بیاہنا} (byāhnā)}{بیاہنا (byāhnā)}]{\ur{بیاہنا} (byāhnā) — to marry off; to wed}

\label{lesson:UR-C488-byahna}



\begin{quote}
Before the new one: say the Urdu for a daughter-in-law, then the Urdu for a son-in-law.
\end{quote}

\subsection*{You'll want to know: \ur{بیاہنا}}
\textbf{\ur{بیاہنا}} — \emph{byāhnā} — \textquotedblleft{}to marry off; to wed\textquotedblright{}.

\begin{quote}
\ur{انہوں} \ur{نے} \ur{بیٹی} \ur{کو} \ur{لاہور} \ur{میں} \ur{بیاہا۔}
\end{quote}

\begin{quote}
\emph{unhoṅ ne beṭī ko lāhaur meṅ byāhā.}
\end{quote}

\begin{quote}
They married their daughter into a family in Lahore.
\end{quote}

Past, with \ur{نے} (he or she …): \ur{اس} \ur{نے} \ur{بیٹے} \ur{کو} \ur{بیاہا}, \emph{us ne beṭe ko byāhā}.

\begin{cousinweb}
\textbf{\ur{بیاہنا}} (\emph{byāhnā}) is inherited from Sanskrit \emph{vivāhayati}, \textquotedblleft{}leads away (a bride), marries\textquotedblright{} (CDIAL 11923).
\end{cousinweb}

\begin{grammarlens}[title={What you've built}]
One more word for describing what you lived and felt.
\end{grammarlens}

\subsection*{Your turn}
\begin{itemize}
\raggedright
  \item \emph{Say it:} \emph{byāhnā}
  \item \emph{Say it:} \emph{byāhnā}, once more
  \item \emph{Say it:} \emph{unhoṅ ne beṭī ko lāhaur meṅ byāhā.} — They married their daughter into a family in Lahore.
  \item \emph{Recall:} say the Urdu for a daughter-in-law, then the Urdu for a son-in-law, then say \emph{byāhnā} again
\end{itemize}

\begin{tcolorbox}[breakable,colback=teal!4,colframe=teal!35!black,title={Before you move on}]
What does \ur{بیاہنا} mean? (\textquotedblleft{}To marry off; to wed\textquotedblright{}.) Say it once more.
\end{tcolorbox}
